\section{Operações em \texorpdfstring{$\mathbb R^n$}{Rn}}

Algumas operações importantes em $\mathbb R^n$ incluem soma, produto por escalar e produto escalar.
Nesta seção, revisaremos tais operações.

\begin{definition}
    Sejam $p=(x_1, \ldots, x_n)$ e $q=(y_1, \ldots, y_n)$ dois elementos de $\mathbb R^n$, e $\alpha \in \mathbb R$ um número real.
    A \emph{soma}\index{soma vetorial} $p+q$ é definida como:
    \begin{equation*}
        p+q =(x_1, \ldots, x_n)+(y_1, \ldots, y_n)=(x_1+y_1, \ldots, x_n+y_n).
    \end{equation*}
    A \emph{multiplicação por escalar}\index{multiplicação por escalar} $\alpha p$ é definida como:
    \begin{equation*}
        \alpha p = \alpha (x_1, \ldots, x_n) = (\alpha x_1, \ldots, \alpha x_n).
    \end{equation*}
    O \emph{produto escalar}\index{produto escalar} $p \cdot q$ é definido como:
    \begin{equation*}
        p \cdot q = (x_1, \ldots, x_n) \cdot (y_1, \ldots, y_n) = x_1y_1 + \ldots + x_ny_n=\sum_{i=1}^n x_iy_i.
    \end{equation*}
\end{definition}

Como usual, escrevemos
\[
    -p=(-1)p
    \quad\text{e}\quad
    p-q=p+(-q).
\]

Desse modo, $-(x_1, \ldots, x_n) = (-x_1, \ldots, -x_n)$ e $(x_1, \ldots, x_n)-(y_1, \ldots, y_n) = (x_1-y_1, \ldots, x_n-y_n)$.

Note que $p+q$ e $\alpha p$ são elementos de $\mathbb R^n$, enquanto $p\cdot q$ é um número real.
Assim, a soma e a multiplicação por escalar produzem vetores, ao passo que o produto escalar associa dois vetores a um escalar.

\begin{example}
    Se $p=(1,-2,3)$, $q=(4,0,-1)$ e $\alpha=-2$, então:
    \[
        p+q=(5,-2,2),
        \qquad
        \alpha p=(-2,4,-6),
    \]
    e
    \[
        p\cdot q
        =1\cdot4+(-2)\cdot0+3\cdot(-1)
        =1.
    \]
\end{example}

\begin{example}
    Se $p=(3,-1,2)$ e $q=(1,4,-2)$, então
    \[
        q-p=(1,4,-2)-(3,-1,2)=(-2,5,-4).
    \]
    Assim, partindo de $p$ e somando o deslocamento $q-p$, chegamos a $q$:
    \[
        p+(q-p)
        =(3,-1,2)+(-2,5,-4)
        =(1,4,-2)
        =q.
    \]
\end{example}

Vamos lembrar de importantes interpretações geométricas da soma e da multiplicação por escalar.

A soma dos vetores $v=(3, 2)$ e $w=(-1, 1)$ pode ser visualizada como o vetor com início na origem e fim no ponto obtido posicionando-se o vetor $w$ com início no ponto final do vetor $v$, conforme ilustrado na Figura~\ref{fig:soma-vetores}.

\subimport{operacoesRn}{somaVetores}
A multiplicação do vetor $v=(1, 2)$ por $2$ e por $-2$ corresponde, respectivamente, a multiplicar o comprimento de $v$ por $2$, mantendo a direção e o sentido, e a multiplicar o comprimento de $v$ por $2$, mantendo a direção e invertendo o sentido, conforme ilustrado na Figura~\ref{fig:produto-por-escalar}.

\subimport{operacoesRn}{produtoPorEscalar}

\begin{example}
    Seja $v=(2,-1)$. Então
    \[
        0v=(0,0),
        \qquad
        \frac12v=(1,-\tfrac12),
        \qquad
        -3v=(-6,3).
    \]
    Multiplicar um vetor por $0$ resulta no vetor nulo.
    
    Multiplicar um vetor por $\frac12$ preserva a direção e o sentido, mas reduz o comprimento pela metade.
    
    Já multiplicar um vetor por $-3$ preserva a direção, inverte o sentido e multiplica o comprimento por $3$.
\end{example}

Agora vamos nos voltar ao produto escalar.

\begin{example}
    O produto escalar entre dois vetores pode ser positivo, nulo ou negativo.
    Por exemplo, tomando $u=(1,2)$, temos:
    \[
        u\cdot(2,1)=1\cdot2+2\cdot1=4>0,
    \]
    \[
        u\cdot(2,-1)=1\cdot2+2\cdot(-1)=0,
    \]
    e
    \[
        u\cdot(-1,-2)=1\cdot(-1)+2\cdot(-2)=-5<0.
    \]
    Por outro lado,
    \[
        u\cdot u=1^2+2^2=5>0.
    \]
\end{example}

Algumas propriedades importantes do produto escalar são as seguintes.

\begin{proposition}\index{produto escalar!propriedades}
    Sejam $p, q, r \in \mathbb R^n$ e $\alpha \in \mathbb R$. Então:
    \begin{enumerate}
        \item $p \cdot q = q \cdot p$ (comutatividade);
        \item $(p+q) \cdot r = p \cdot r + q \cdot r$ (distributividade);
        \item $(\alpha p) \cdot q = \alpha(p \cdot q)$ (homogeneidade em relação a escalares);
        \item $p \cdot p \geq 0$.
    \end{enumerate}
\end{proposition}

\begin{proof}
    Escrevamos as coordenadas de $p, q, r$ como $p=(x_1, \ldots, x_n)$, $q=(y_1, \ldots, y_n)$ e $r=(z_1, \ldots, z_n)$.

    Para a comutatividade, temos:
    \begin{equation*}
        p \cdot q = \sum_{i=1}^n x_iy_i = \sum_{i=1}^n y_ix_i = q \cdot p.
    \end{equation*}

    Para a distributividade, temos:
    \begin{align*}
        (p+q) \cdot r &= \sum_{i=1}^n (x_i+y_i)z_i \\
        &= \sum_{i=1}^n x_iz_i + \sum_{i=1}^n y_iz_i \\
        &= p \cdot r + q \cdot r.
    \end{align*}

    Quanto à homogeneidade em relação a escalares, temos:
    \begin{equation*}
        (\alpha p) \cdot q = \sum_{i=1}^n (\alpha x_i)y_i = \alpha \sum_{i=1}^n x_iy_i = \alpha (p \cdot q).
    \end{equation*}

    Finalmente, como $x_i^2\geq 0$ para todo $i$, temos:
    \begin{equation*}
        p \cdot p = \sum_{i=1}^n x_i^2 \geq 0.
    \end{equation*}
\end{proof}

\subsection*{Exercícios}

\begin{exercise}
    Sejam $p=(2,-1,0,3)$ e $q=(-1,4,5,2)$ elementos de $\mathbb R^4$.
    Calcule:
    \begin{enumerate}[label=(\alph*)]
        \item $p+q$;
        \item $p-q$;
        \item $q-p$;
        \item $-3p$;
        \item $p\cdot q$;
        \item $q\cdot q$.
    \end{enumerate}
\end{exercise}

\begin{exercise}
    Sejam $p=(x_1,\ldots,x_n)$ e $q=(y_1,\ldots,y_n)$ elementos de $\mathbb R^n$.
    Mostre que
    \begin{enumerate}[label=(\alph*)]
        \item $p+(q-p)=q$;
        \item $q-p=-(p-q)$.
    \end{enumerate}
\end{exercise}

\begin{exercise}
    Seja $u=(1,2)$.
    Determine todos os vetores $v=(x,y)\in\mathbb R^2$ tais que $u\cdot v=0$.
\end{exercise}

\begin{exercise}
    Sejam $p,q\in\mathbb R^n$.
    Mostre que:
    \[
        (p+q)\cdot(p+q)=p\cdot p+2p\cdot q+q\cdot q.
    \]
\end{exercise}

\begin{exercise}
    Sejam $p,q\in\mathbb R^n$.
    Mostre que:
    \[
        (p-q)\cdot(p-q)=p\cdot p-2p\cdot q+q\cdot q.
    \]
\end{exercise}

\begin{exercise}
    Seja $p\in\mathbb R^n$.
    Mostre que $p\cdot p=0$ se, e somente se, $p=0$.
\end{exercise}

\begin{exercise}
    Para cada $i\in\{1,\ldots,n\}$, seja $e_i\in\mathbb R^n$ o vetor cuja $i$-ésima coordenada é $1$ e cujas demais coordenadas são $0$.
    Mostre que:
    \[
        e_i\cdot e_j=
        \begin{cases}
            1, & \text{se } i=j,\\
            0, & \text{se } i\neq j.
        \end{cases}
    \]
    Em seguida, mostre que, se $p=(x_1,\ldots,x_n)$, então
    \[
        p\cdot e_i=x_i.
    \]
\end{exercise}
